\documentclass[12pt]{article}
\usepackage[margin=1in]{geometry}
\usepackage{amsmath}
\usepackage{amssymb}

\begin{document}

\section*{ECO 310 - Problem Set 2}
\subsection*{Question 5(a) - Answer}

\textbf{Problem:} Consider a decision-maker who expects to have a car accident with chance $\pi$; if this occurs, he will incur \$L in damages. He can purchase as much auto insurance, $q$, as he likes, at a price of $p$ per dollar of coverage: this means that if he pays $pq$ upfront (as the insurance premium), he'll receive a payment of $q$ from the insurance company if an accident occurs.

\vspace{0.3cm}

\textbf{Question:} Write out his expected utility from purchasing insurance level $q$, assuming a utility-of-wealth function $u(w)$ and initial wealth $w_0$.

\vspace{0.5cm}

\textbf{Answer:}

There are two possible states of the world:

\begin{enumerate}
    \item \textbf{No accident} (probability $1-\pi$):
    \begin{itemize}
        \item The decision-maker pays the insurance premium $pq$
        \item Final wealth: $w_0 - pq$
        \item Utility: $u(w_0 - pq)$
    \end{itemize}

    \item \textbf{Accident occurs} (probability $\pi$):
    \begin{itemize}
        \item The decision-maker pays the insurance premium $pq$
        \item Incurs damages of $L$
        \item Receives insurance payout of $q$
        \item Final wealth: $w_0 - pq - L + q = w_0 - L + q(1-p)$
        \item Utility: $u(w_0 - L + q(1-p))$
    \end{itemize}
\end{enumerate}

\vspace{0.3cm}

Therefore, the expected utility from purchasing insurance level $q$ is:

\begin{equation}
\boxed{EU(q) = (1-\pi)u(w_0 - pq) + \pi u(w_0 - L + q(1-p))}
\end{equation}

or equivalently:

\begin{equation}
\boxed{EU(q) = (1-\pi)u(w_0 - pq) + \pi u(w_0 - L - pq + q)}
\end{equation}

\end{document}
